\section{The Solar and the Demetrian Races}

\begin{quotex}
\textbf{Julius Evola} begins his discussion of the spiritual races in his \emph{Sintesi di dottrina della razza} with the Solar and Demetrian races. He gives and almost Heideggerian take on the \emph{avidya} teaching of the Vedanta. To recast it in these terms: The Hyperborean experiences a feeling of homelessness, arising from the ignorance or forgetfulness of Being, and finds himself thrown among beings.

This homeland, from the North, should be understand in a metaphysical, rather than a temporal, sense. Evola repeats the error we pointed out about the ritual between the priest and the king as described in the Brahmanas. Certainly the priesthood can decay (or has decayed) into a lunar type, but are (or were) there not sages or wizards with the solar powers described?
\end{quotex}

In order to move onto something more specific, we will say that the doctrine of the third degrees of race must essentially limit its researches to the sphere of influence of a given race of the spirit and of its primordial tradition, following its developments, the mutations (paravariations) and also the distortions in the cycle that coincides with them and in which they act and react in the confrontation with the influences of different races or new environmental conditions. Once research is circumscribed that way, we arrive at a more limited concept of race, corresponding to that of the various differentiations or articulations of the primary element of such a cycle. It is natural that likewise we cannot think of an atomic separation of the various ``races of the spirit'': their differences are not such as to exclude relationships not only of derivation, but even of different hierarchical rank.

We already outlined the typology of races of the spirit, as far as it concerns the specific human cycle of the hyperborean race, in the second part of our work \emph{Revolt against the Modern World} (with special regard to the properly traditional and spiritual aspect), both in the choice of the writings of J. J. Bachofen, and in the related interpretation of them in a racial sense, as understood in the previously mentioned volume we edited under the title \emph{La razza solare}. For a fuller treatment, the reader is therefore referred back to these two works. Here we will give only a brief schematic synthesis, necessarily lacking any justifying elements.

The solar or Olympic-solar race, corresponding to the hyperborean lineage and tradition, is to be considered superior and anterior to all the others in the cycle in question. It has as characteristic a type of ``natural supernaturality''; spirit and power, calm dominance, and readiness for precise and absolute action, a sense of ``centrality'' and of ``unshakeability'', and in its exterior effects, that virtue to which the ancients attributed a ``numinous'' quality (from \emph{numen}), as superiority that imposes itself directly and irresistibly, that awakens simultaneously both terror and veneration, are marks of this ``race of the spirit'', by means of which it is naturally predestined to command and, at the limit, to the regal function. Fire and ice are united in it as in the vague symbols of the original Nordic homeland and of the cycle in which it is eminently and primordially manifested: ice, as transcendence and inaccessibility; fire, as the radiant specifically solar quality of the beings who give, who awaken life, and generate light, but always from a supreme distance and almost with indifference, as in the wake of a boat, not through zeal, inclination, or human concern.

The ancient symbol of gold always had connections with this form of spirituality. In the political forms of the origins, it acted as the substrate to sacred, or divine, regality, i.e., to the union of the two powers, of the regal and priestly function, the latter understood in a higher sense that will soon be made clear. The symbolic designations of the ``divine'' or ``celestial'' race for it must refer to the absence of the dualistic sentiment in relations to supernatural reality, something that however is quite distinct from everything that, in the modern sense, is immanence or Promethean velleity: it is not about men who believe they are gods, but of types who naturally, through a still not dim memory of origins and a condition of soul and body that do not cripple such a memory, feel they do not properly belong to the terrestrial race, to the extent of believing themselves men only by accident, or through ``ignorance'', or ``sleep''. The two terms \emph{vidya} and \emph{avidya} of the ancient Indo-Aryan teaching, that mean respectively ``knowledge'' (of the ``supreme identity'') and ``ignorance'' (that leads to the self-identification with one of the forms or modes of being of the conditioned world), are to be understood exactly in this connection: when these are applied to a different human condition and a different race of the spirit, that is, when turned into ``philosophical'' terms, they lose every meaning and give rise to misunderstandings of various types.

The most recent ``spiritual races'' of the cycle, to which our contemporaries also belong, have as a prerequisite a division and a separation of the two elements of ``spirituality'' and ``virility''---and also ``transcendence'' and ``humanity''---that are synthetically reunited in the solar race. First of all, let's specify the lunar race or demterian race. According to the analogous relation, while the solar element is that which has in itself its own light and, in general, its own principle, the sun being, in such regard, the center of a given planetary system, the lunar element is that which instead receives or draws from another its light and principle. In the lunar race the sense of spiritual centrality is therefore lost, or through degeneration (the moon as the sun extinguished), or through passive crossbreeding with races of other cycles, of the ``telluric'' type, that have degraded its originary solar quality. The moon, as Bachofen points out, was also called the ``celestial earth'' by the ancients. There is therefore a sublimation of the law of the earth, the fate that is disclosed under the types of cosmic harmony and natural law: man here no longer feels himself the active center of spiritual reality: he is not this same reality, but rather the one who contemplates it, who studies its laws, who surpasses material action with contemplation and ``tellurism'', but still does not reach spiritual action. The adjective ``demetrian'', that we likewise give to this race, refers to a spirituality of a diffused, pantheistic, less dominating character that is pervaded by the sense of cosmic naturalistic laws and by a sacredness placed essentially under the feminine sign: spirituality that was precisely typical of the ancient demetrian cults. By extension, the priestly man is lunar in opposition to the regal man, he is the man who in the face of the spirit behaves as a normal woman facing a man, i.e., with a sense of submission and devotion. It is then interesting to note that the ancient traditions associated what today would be called cerebrality or intellectuality with the moon, labeling instead the heart with the sun and referring higher form of consciousness to this. The lunar type is in fact also intellectual, the man of passive reflection who, as the word says it, moves only among ``reflections'', among shadows of ideas and of things.

Therefore the characteristics of the lunar race are varied. In the political camp, wherever a division between the temporal and priestly powers is perceived, the lunar spirit inevitably arises. The lunar is the dominator who receives the supreme consecration of its power from the other, from a distinct priestly caste, which it does not give it to itself. In general, the lunar man has spiritually feminine traits. He lacks the feeling of centrality. In correspondence to the races of the body, he has the predominant nature of the demetrian race which is of the stock we called Atlantic-Occidental, in its prehistoric forms that lead us, e.g., up to the Pelasgic, Minoan-Mycenaean, or Etruscan civilization and to its last revival, among which Pythagoreanism returns. Such a race represents a decay of the hyperborean spirituality which appeared suddenly in the regions of Atlantis through processes of action and reaction, causing a series of other changes. Lunar elements can moreover be observed in the race of the man of the East (Alpine-Oriental) by some racialists---the psycho-anthropology of Claus designates this race as the ``race of evasion'' --- \emph{der Enthebungsmenschen} --- something that corresponds visibly to a characteristic of the lunar man.


\flrightit{Posted on 2014-03-26 by Aeneas}

